% ── Section V: Results ───────────────────────────────────────────────
% Length  : ~1.5 pages
% Fill AFTER running experiments in Month 4.
% Structure:
%   5.1  Ablation Study Results    (main table)
%   5.2  Comparison with Prior Systems
%   5.3  Qualitative Results       (figures)
%   5.4  Runtime Analysis
% ─────────────────────────────────────────────────────────────────────

\section{Results}
\label{sec:results}

% ── 5.1 ─────────────────────────────────────────────────────────────
\subsection{Ablation Study}
\label{subsec:ablation_results}

Table~\ref{tab:ablation_results} reports the quantitative results
of the four ablation configurations on the VN-MixedTraffic test set.

\begin{table}[h]
  \centering
  \caption{Ablation results on VN-MixedTraffic test set.
           \textbf{Bold} = best. $\uparrow$ higher is better,
           $\downarrow$ lower is better.}
  \label{tab:ablation_results}
  \begin{tabular}{lcccc}
    \toprule
    Config & MOTA$\uparrow$ & IDF1$\uparrow$ & Acc-F1$\uparrow$ & FP Rate$\downarrow$ \\
    \midrule
    Baseline (Exp~0) & -- & -- & -- & -- \\
    Exp~1 (+C1)      & -- & -- & -- & -- \\
    Exp~2 (+C1,C2)   & -- & -- & -- & -- \\
    Exp~3 (Proposed) & \textbf{--} & \textbf{--} & \textbf{--} & \textbf{--} \\
    \bottomrule
  \end{tabular}
\end{table}

% TODO: analysis paragraphs
%   - What does adding C1 alone buy?  (FP rate should drop most here)
%   - What does adding C2 add?        (MOTA and IDF1 should improve most here)
%   - What does adding C3 complete?   (Acc-F1 should improve most here)


% ── 5.2 ─────────────────────────────────────────────────────────────
\subsection{Comparison with Prior Systems}
\label{subsec:comparison}

\begin{table}[h]
  \centering
  \caption{Comparison with prior methods.}
  \label{tab:comparison}
  \begin{tabular}{lccccc}
    \toprule
    Method & Mixed & Class-cond. & Adaptive & Severity & MOTA$\uparrow$ \\
    \midrule
    Hikvision AID~\cite{Hikvision_AID}  & --         & --         & --         & --         & N/A \\
    ByteTrack~\cite{Zhang2022_bytetrack} & \checkmark & --         & --         & --         & --  \\
    YOLO-only                            & \checkmark & --         & --         & --         & --  \\
    Proposed                             & \checkmark & \checkmark & \checkmark & \checkmark & \textbf{--} \\
    \bottomrule
  \end{tabular}
\end{table}

% TODO: narrative — what each comparison reveals


% ── 5.3 ─────────────────────────────────────────────────────────────
\subsection{Qualitative Results}
\label{subsec:qualitative}

% TODO: include 2 figures
%   fig5_qualitative_tracking.pdf — side-by-side: baseline vs proposed,
%                                   track IDs over weaving sequence
%   fig6_failure_cases.pdf        — motivation: annotated frames showing
%                                   baseline FP and track loss

Fig.~\ref{fig:qualitative} illustrates representative tracking sequences
comparing the baseline and the proposed method.
% \begin{figure}[t]
%   \centering
%   \includegraphics[width=\linewidth]{figures/fig5_qualitative_tracking.pdf}
%   \caption{Tracking comparison on a dense motorcycle scene.
%            (a)~Baseline: track fragmentation visible (ID switches).
%            (b)~Proposed: stable IDs maintained across weaving.}
%   \label{fig:qualitative}
% \end{figure}


% ── 5.4 ─────────────────────────────────────────────────────────────
\subsection{Runtime Analysis}
\label{subsec:runtime}

\begin{table}[h]
  \centering
  \caption{Runtime analysis (TODO: hardware spec).}
  \label{tab:runtime}
  \begin{tabular}{lcc}
    \toprule
    Component & Latency (ms) & \% of total \\
    \midrule
    YOLOv8 inference    & -- & -- \\
    Context Estimator   & -- & -- \\
    Modified ByteTrack  & -- & -- \\
    Reliability Scorer  & -- & -- \\
    EfficientNet-B0     & -- & -- \\
    \midrule
    \textbf{Total}      & \textbf{--} & 100\% \\
    \bottomrule
  \end{tabular}
\end{table}

% TODO: note real-time threshold (30 FPS = 33 ms/frame)
%       confirm pipeline meets real-time requirement
